\section{Boundary waves for the rotation embeddings}\label{app:local}
This appendix proves \cref{lem:passages}. The candidate arrays below are explicit formulas checked by the actual cutoff rule \eqref{eq:row-rule}; composing their array substitutions is not an assumption that overlapping rowmotion waves commute.

\subsection{An isolated passage with a passive close pair}\label{app:ordinary}
Suppose $a_i,b_i$ are strict decreasing valid endpoint lists with $b_1=n$, and let $k$ satisfy
\[
 a_i\le b_{i+1}\quad(i<k),\qquad
 a_k\ge b_{k+1}+3\quad(k<m).
\]
Assume $a_m\ge1$, $b_2+m<n$ (when $m\ge2$), and $a_k+m-2<n$. For $0\le t<k$, define
\begin{equation}\label{eq:ordinary-upper}
 A_i(t)=a_i+t,\qquad
 B_i(t)=\begin{cases}b_{i+1}+t&i\le t,\\b_i+t&i>t.\end{cases}
\end{equation}
For $k\le t<m$, define
\begin{align}
 A_i(t)&=\begin{cases}a_{i+1}+t&k\le i\le t,\\a_i+t&\text{otherwise},\end{cases}\label{eq:ordinary-lower}\\
 B_i(t)&=\begin{cases}b_{i+1}+t&i<k,\\a_k+t-2&i=k,\\b_i+t&i>k.\end{cases}\notag
\end{align}
The terminal state is
\begin{align}
 A_i(m)&=\begin{cases}a_i+m&i<k,\\a_{i+1}+m&k\le i<m,\\1&i=m,\end{cases}\label{eq:ordinary-terminal}\\
 B_i(m)&=\begin{cases}b_{i+1}+m&i<k,\\a_k+m-2&i=k,\\b_i+m&i>k.\end{cases}\notag
\end{align}

Here is the full cutoff induction. Before the switch, only ceiling row $h=t+1$ is missing. For $h<k$, its next ceiling gives $D_h=b_{h+1}+t+1>A_h$, and its minimum qualifies. Row $h$ becomes $[a_h+t+1,b_{h+1}+t+1]$, while all other rows translate. At $h=k<m$, $A_h\ge D_h+2$, so that minimum fails to qualify, and $E_h=a_{h+1}+t$. The incomparable left interval joins the part below the ceiling, giving $[a_{k+1}+t+1,a_k+t-1]$. At $k=m$, the zero cutoffs give $[1,a_m+m-2]$. After switching, every ceiling row is present and only minimum row $t+1$ is missing. Its $E$ cutoff reads the next lower endpoint, or $0$ at the bottom. This propagates the lower equality and gives \eqref{eq:ordinary-terminal}.

All rows are valid: before switching use $a_i\le b_{i+1}$; at the switch use $a_{k+1}\le b_{k+1}\le a_k-3$; afterwards each lower source moves downward. The flipped upper endpoint is strictly between its neighbors because $b_k\ge a_k>a_k-2>b_{k+1}$. All remaining list comparisons are inherited. The strict top bounds exclude new ceiling defects. Thus the formulas allow singleton rows and arbitrarily small positive same-list gaps.

For an isolated matching passage, the last closing $N$ is matched to the first balanced-suffix opening $p_k$. The stack rule gives $p_i<q_{i+1}$ for $i<k$, and $p_k>q_{k+1}$ when $k<m$. An opening and closing at distance one would be an excluded adjacent matched peak, so decoding yields $a_i\le b_{i+1}$ before the switch and $a_k\ge b_{k+1}+3$ at it. Large neighboring support gaps give the strict top bounds and $a_m\ge D-m+1$. Under $\rho^m$, $N$ becomes the new opening $m$, while $p_k$ changes from opening to closing. The two decoder offsets differ by two, so the resulting endpoints are precisely \eqref{eq:ordinary-terminal}. This proves \cref{lem:passages}(i), including every passive location of the close pair.

\subsection{A short matched leaf}\label{app:leaf}
Suppose $m\ge2$, $\delta\ge2$, $H\ge m+6$, and
\[
 b_1=n,\qquad a_1=n-\delta+2,\qquad a_m\ge H.
\]
All other endpoints are at least $H$ below $a_1$, and pairwise separated by at least $H$. Set $T=m+\delta$. At time $0$ use the input. For $1\le t<T$, define
\begin{equation}\label{eq:leaf-A}
 A_i(t)=\begin{cases}
 a_{i+1}+t&i\le t<m,\\
 a_i+t&t<m,\ i>t,\\
 a_{i+1}+t&t\ge m,\ i<m,\\
 t-m+1&t\ge m,\ i=m.
 \end{cases}
\end{equation}
Put $u=\max(0,t-\delta)$, and define
\begin{equation}\label{eq:leaf-B}
 B_i(t)=\begin{cases}
 b_{i+1}+t&i\le u,\\
 n-\delta+t&u=0,\ i=1,\\
 b_i+t&\text{otherwise}.
 \end{cases}
\end{equation}
Then the terminal state is
\begin{equation}\label{eq:leaf-terminal}
 (A_i(T),B_i(T))=
 \begin{cases}(a_{i+1}+T,b_{i+1}+T)&i<m,\\(1,\delta-1)&i=m.
 \end{cases}
\end{equation}

At time $0$ the only missing ceiling is row $1$, with $D_1=b_2+1<a_1$ and $E_1=a_2$. Its output is $[a_2+1,a_1-1]$, giving the first lower equality and \eqref{eq:leaf-B}. For $1\le t<m$, the missing minimum is row $l=t+1$. Its next qualifying minimum has source $a_{l+1}+t$, or the zero sentinel if $l=m$. Its ceiling is present, so the lower equality advances as in \eqref{eq:leaf-A}.

Before $t=\delta$, all ceilings are present. At $t=\delta$ the first upper endpoint reaches $n$; thereafter the missing ceiling is row $h=u+1$. Until the last step, $h<m$ and
\[
 A_h(t)=a_{h+1}+t,\qquad D_h=b_{h+1}+t+1,\qquad E_h=A_h(t).
\]
The possible missing minimum lies $\delta\ge2$ rows lower, so the displayed minimum is present. As $a_{h+1}\le b_{h+1}$, the incomparable contribution is removed. Thus the upper equality advances. At the final input $t=T-1$, row $m$ has $A_m=\delta$ and missing ceiling, with $D_m=E_m=0$. Its output is the complement of $[\delta,B_m]$ in $[1,B_m]$, namely $[1,\delta-1]$.

These are all exceptional cutoffs. Every changed lower source comes from a later original row, or is the inserted $t-m+1$. The top-row flip is valid since $a_1-a_2\ge H$. Retained upper endpoints satisfy $b_2+T\le n+m+2-H<n$, and the special $n-\delta+t$ reaches $n$ only at its intended second passage. This proves validity and the full induction.

If the close pair in the matching is matched, it is the last $AB$ leaf. After rotating by $T$, its endpoints are $m,m+\delta$, so its decoded row is $[1,\delta-1]$, while all other endpoints translate. This is \eqref{eq:leaf-terminal}. For $\delta\ge m$, the formulas on $m\le t\le\delta$ give
\[
 A_i=a_{i+1}+t\ (i<m),\quad A_m=t-m+1,\quad
 B_1=n-\delta+t,\quad B_i=b_i+t\ (i>1),
\]
which is exactly the intermediate direct decoder. This proves parts (ii)--(iii) for the matched case.

\subsection{Two close closings from distinct edges}\label{app:distinct}
Now let $m\ge2$, $\delta\ge1$, and $H\ge m+6$. Start from valid strictly decreasing endpoint lists $a_i,b_i$ satisfying
\[
 b_1=n,\qquad b_2=n-\delta,\qquad a_m\ge H.
\]
Every other endpoint lies at most at $b_2-H$, and every pair other than $b_1,b_2$ is separated by at least $H$. Let
\[
 k=\min\{i:i=m\text{ or }a_i>b_{i+1}\},\qquad
 r=\min\{i:1\le i<k,\ a_i>b_{i+2}\},
\]
where $b_{m+1}=0$ is used only in defining $r$. Then $2\le k\le m$ and $1\le r<k$: the last eligible $i=k-1$ always satisfies the inequality. Put $T=m+\delta$.

\paragraph{Candidate phases.}
Subtract the common translation $t$ and first define older-wave arrays:
\begin{equation}\label{eq:older-x}
 x_i(t)=\begin{cases}
 a_i&t<k,\\
 a_{i+1}&k\le t<m,\ k\le i\le t,\\
 a_i&k\le t<m,\ i\notin[k,t],\\
 a_i&t\ge m,\ i<k,\\
 a_{i+1}&t\ge m,\ k\le i<m,\\
 1-m&t\ge m,\ i=m,
 \end{cases}
\end{equation}
\begin{equation}\label{eq:older-y}
 y_i(t)=\begin{cases}
 b_{i+1}&t<k,\ i\le t,\\
 b_i&t<k,\ i>t,\\
 b_{i+1}&t\ge k,\ i<k,\\
 a_k-2&t\ge k,\ i=k,\\
 b_i&t\ge k,\ i>k.
 \end{cases}
\end{equation}
For $0\le t<T$ put $u=t-\delta$ and set
\begin{equation}\label{eq:distinct-A}
 A_i(t)=\begin{cases}x_{i+1}(t)+t&r\le u,\ r\le i\le u,\\x_i(t)+t&\text{otherwise},\end{cases}
\end{equation}
\begin{equation}\label{eq:distinct-B}
 B_i(t)=\begin{cases}
 y_i(t)+t&u<0,\\
 y_{i+1}(t)+t&0\le u<r,\ i\le u,\\
 y_i(t)+t&0\le u<r,\ i>u,\\
 y_{i+1}(t)+t&u\ge r,\ i<r,\\
 a_r+t-2&u\ge r,\ i=r,\\
 y_i(t)+t&u\ge r,\ i>r.
 \end{cases}
\end{equation}
Since $u<m$ here, no undefined lower source is read. At $t=T$ use
\begin{align}
 A_i(T)&=\begin{cases}
 a_i+T&i<r,\\
 a_{i+1}+T&r\le i<k-1,\\
 a_{i+2}+T&k-1\le i\le m-2,\\
 \delta+1&i=m-1,\\1&i=m,
 \end{cases}\label{eq:distinct-terminal-A}\\
 B_i(T)&=\begin{cases}
 b_{i+2}+T&i<r,\\a_r+T-2&i=r,\\b_{i+1}+T&r<i<k,\\
 a_k+T-2&i=k,\\b_i+T&i>k.
 \end{cases}\label{eq:distinct-terminal-B}
\end{align}
Empty index ranges are omitted. These formulas include $m=2$, $r=1$, $k=2$, $\delta=1$.

\paragraph{The inequalities used in the verification.}
The hypotheses concern the \emph{original} endpoints. In particular they give
\begin{align*}
 a_i+H&\le b_i &&(1\le i\le m),&
 a_{i+1}+H&\le a_i &&(i<m),\\
 a_i+H&\le b_{i+1} &&(i<k),&
 b_{k+1}+H&\le a_k &&(k<m),\\
 a_i+H&\le b_{i+2} &&(i<r),&
 b_{r+2}+H&\le a_r &&(r+2\le m).
\end{align*}
Also $a_r+H\le b_{r+1}$ because $r<k$, and $a_k+H\le a_r$ because $r<k$. These inequalities follow by combining the definitions of the first crossings with pairwise separation; none applies the separation hypothesis to a shifted endpoint. In particular,
\[
 (a_k-2)-b_{k+1}\ge H-2\quad(k<m),\qquad
 (a_r-2)-b_{r+2}\ge H-2\quad(r+2\le m).
\]
The lower bound after a flip is $H-2$, not $H$. It is positive since $H\ge m+6$. The complementary comparisons are
\[
 b_k-(a_k-2)\ge H+2,\quad
 b_{r+1}-(a_r-2)\ge H+2,\quad
 (a_r-2)-(a_k-2)\ge H.
\]
They keep each flipped upper source between its neighbors, including the case $k=r+1$ in which the two flipped sources are neighbors.

\paragraph{Validity of the candidate rows.}
Write $L_i=x_i(t)+t$ and $V_i=y_i(t)+t$ for the older arrays. For $i<k$, a changed upper source is $b_{i+1}$ and $a_i+H\le b_{i+1}$. At row $k$ after switching, the lower source is at most $a_{k+1}$ if $k<m$, so $a_{k+1}\le b_{k+1}\le a_k-H<a_k-2$. If $k=m$, its new lower endpoint is $t-m+1<a_m+t-2$, since $a_m\ge H$. The remaining rows retain their upper sources and only move lower sources to later indices. Thus $L_i\le V_i$, and the older lower array is nonincreasing with $L_i\le a_i+t$.

The younger operation has three row types. Where it shifts an upper source and $i<r$, we have $L_i=a_i+t$ and $V_{i+1}=b_{i+2}+t$, because $i+1<k$ and $i+1\le t$. Hence $L_i+H\le V_{i+1}$. At its flipped row $r$, the lower source is $L_{r+1}\le a_{r+1}+t\le a_r+t-H<a_r+t-2$. For $i>r$, its lower source is either $L_i$ or $L_{i+1}\le L_i$, while its upper source remains $V_i$. Before the younger switch only the first of these upper-shift cases occurs. This proves row nonemptiness throughout \eqref{eq:distinct-A}--\eqref{eq:distinct-B}. The deletion of entries preserves lower order; the comparisons in the preceding paragraph prove upper order at the two flipped entries, and all other comparisons are inherited or deliberate repetitions.

Every old lower source is at least $H$. The inserted lower value $t-m+1$ is used only for $t\ge m$, so it is positive and is smaller than every surviving old lower source at the same time. The source $b_1$ occurs only at $t=0$; $b_2$ occurs only through $t=\delta$, when $b_2+t=n$. Every other original source is at most $n-\delta-H$, so even at $T=m+\delta$ its translate is at most $n+m-H<n$. A flipped source is two smaller still. These bounds place the arrays in $[1,n]$. At $T$, the new bottom lower values satisfy $1<\delta+1<a_i+T$ for every surviving old lower source, and their upper entries are at least $a_m+T-2>\delta+1$. Thus the terminal arrays are valid as well.

\paragraph{Missing rows.}
The strict original comparisons, with the positive residual gap $H-2$ at flips, leave exactly the displayed repetitions. They give the following complete list:
\begin{center}
\begin{tabular}{lll}
\toprule
Front & Missing ceiling & Missing minimum\\
\midrule
Older & row $t+1$ if $0\le t<k$ & row $t+1$ if $k\le t<m$\\
Younger & row $u+1$ if $0\le u<r$ & row $u+1$ if $r\le u<m$\\
\bottomrule
\end{tabular}
\end{center}
A front that has exited contributes no missing row. The front positions differ by $\delta\ge1$, so a row is never both kinds of defect. For $\delta=1$, two same-kind missing rows can be adjacent. The associated three equal endpoints form a single plateau; the cutoff must skip both missing rows.

\paragraph{Ceiling cutoffs and qualification, including the switches.}
A row with a ceiling has $D_i=U_i=B_i+1$. In a missing-ceiling row, $U_i-1=B_i$, because the upper endpoint repeats the preceding one or belongs to the top plateau at $n$. Its $D_i$ is the next present ceiling below the entire plateau, or zero if there is none.

At the older upper front $h=t+1\le k$, the younger front is above $h$, so
\[
 A_h=a_h+t,\qquad
 D_h=\begin{cases}b_{h+1}+t+1&h<m,\\0&h=m.\end{cases}
\]
For $h<k$, $D_h-A_h\ge H+1$, and the present minimum qualifies for $\Delta C$. For $h=k<m$, $A_h-D_h\ge H-1>0$, so it fails. For $h=k=m$, $A_h\ge H+t>0=D_h$, so it also fails. This accounts for both the ordinary and bottom older switch.

At the younger upper front $h=u+1\le r$, the lower source is $A_h=a_h+t$. The complete cutoff cases are
\[
 D_h=\begin{cases}
 a_k+t-1&t\ge k\text{ and }h=k-1,\\
 b_{h+2}+t+1&\text{otherwise, if }h+2\le m,\\
 0&\text{otherwise, if }h+2>m.
 \end{cases}
\]
In the first case the next present upper source is the already flipped $a_k-2$. In the second case the next source is $b_{h+2}$, including when adjacent ceiling holes require skipping two rows. The third case is the bottom sentinel: it is zero, not $b_{m+1}+t+1$.

If $h<r$, only the second case can occur, since $h+2\le r+1\le k\le m$ and $h\ne k-1$. It gives $D_h-A_h\ge H+1$. At $h=r$ the first case gives $A_h-D_h=a_r-a_k+1\ge H+1$; the second gives $A_h-D_h\ge H-1$; the third gives $A_h>0=D_h$. Thus the younger minimum qualifies exactly for $h<r$ and fails exactly at $h=r$. In particular, when $\delta=1$ and $r=k-1$, the simultaneous switching input is $t=k-1$, $u=r-1$. The older wave has \emph{not} yet switched at that input. Both upper holes skip to the same next ceiling $b_{k+1}+t+1$ if $k<m$, and to zero if $k=m$; both minima fail qualification.

\paragraph{Lower cutoffs and the one-step recurrence.}
Present minima in rows with a ceiling qualify because that ceiling has cutoff $B_i+1>A_i$. Present minima at nonswitching upper holes qualify by the older- and younger-front inequalities just proved. Thus the only failed present minima are the switching rows. Define the set of failed or missing minimum rows at time $t<T$ by
\[
 Q_t=\{t+1:k-1\le t<m\}\ \cup\
     \{u+1:0\le u<m,\ r-1\le u\}.
\]
A set term is omitted when its stated condition fails. For each row $i$, let $s(i)$ be the least index $j\ge i$ outside $Q_t$, if one exists. The minimum cutoff is exactly
\[
 E_i=\begin{cases}A_{s(i)}&s(i)\text{ exists},\\0&\text{otherwise}.
 \end{cases}
\]
For a row outside $Q_t$ this reads $A_i$. A single hole reads $A_{i+1}$, unless it is the bottom row, when it reads zero. For two adjacent holes $i,i+1$, both read $A_{i+2}$, or zero if $i+1=m$. Two separated holes each read the next row. This exhausts the possibilities because $Q_t$ has at most two elements.

For any missing-ceiling row the incomparable remainder in \eqref{eq:row-rule} is $[D_i+1,A_i-1]$ when this interval is nonempty. It joins $[E_i+1,D_i]$. For a row with a ceiling the incomparable part is empty. Therefore in both cases the output endpoints are
\begin{equation}\label{eq:distinct-recurrence}
 A_i(t+1)=E_i+1,\qquad
 B_i(t+1)=\max(D_i,A_i-1).
\end{equation}
Here are the source substitutions in this recurrence. Before an upper switch its minimum qualifies, so the lower endpoint translates. At a switch or existing lower hole, the preceding scan deletes the old lower source and copies the next surviving source. A single hole consequently advances one lower deletion in \eqref{eq:distinct-A}; adjacent holes advance both deletions and produce the same copied value in their two rows. A bottom hole inserts $1$. For $\delta=1$, the two holes at input $t=m-1$ are rows $m-1,m$ and both read zero, producing lower endpoints $1,1$ at time $m$. At that final input only row $m$ is a hole; the next last two lower endpoints are therefore $2,1$, as \eqref{eq:distinct-terminal-A} requires. If $\delta>1$, the older bottom value at the final input is $\delta$ and translates to $\delta+1$ in row $m-1$, while the younger bottom hole inserts $1$ in row $m$. These are the two terminal lower cases; no source $a_{m+1}$ is used.

For upper endpoints, a present ceiling gives $B_i(t+1)=B_i+1$. A nonswitching upper hole has $A_i<D_i$, so its output is the next present upper source plus one, advancing the corresponding deletion in \eqref{eq:distinct-B}. At the older or younger switch, $A_i>D_i$ and the output is $A_i-1$. Relative to translation by $t+1$, its source is respectively $a_k-2$ or $a_r-2$. Adjacent nonswitching upper holes copy the same ceiling, whereas simultaneous switches use their distinct old $A_i$ and the shared lower cutoff. These substitutions give \eqref{eq:distinct-A}--\eqref{eq:distinct-B} at time $t+1<T$, and \eqref{eq:distinct-terminal-A}--\eqref{eq:distinct-terminal-B} at $t+1=T$. Together with the initial arrays they prove the full induction for the actual action.

\paragraph{Rerooting interpretation.}
If the two largest support points belong to distinct edges, their last letters are $BB$: a last $AB$ would be a matched leaf. Their edges are nested. The first balanced-suffix scan identifies $k$ as the mate rank of the last closing. Removing that outer pair and scanning again identifies $r$ as the mate rank of the penultimate closing, giving exactly its defining comparison above. At time $T$, only these two closings have crossed the cut. They become openings at $m+\delta,m$, giving lower endpoints $\delta+1,1$; their former mates become closings, giving upper endpoints $a_k+T-2,a_r+T-2$. All other endpoint labels translate. This is \eqref{eq:distinct-terminal-A}--\eqref{eq:distinct-terminal-B}.

For $\delta\ge m$ and $m\le t\le\delta$, $u\le0$ and the older wave has finished. The displayed arrays reduce to its ordinary first-reroot decoder translated from time $m$. This includes $\delta=m$ and both endpoints of the intermediate range. Thus the macro agrees at every clear cut, proving the remaining cases of \cref{lem:passages}.
